\section{Trust Boundary and Adversary Model}
\label{sec:threat}

\silr{} relocates the trust boundary from the LLM to the shadow simulator: no admission verdict depends on the LLM's reasoning, only on the simulator's deterministic response. We fix the adversary here---before the mechanism (Sec.~\ref{sec:method})---since it shapes what the gate must defend. The LLM is untrusted; the simulator, verifier, checkers, and tool layer are trusted.

\paragraph{Adversary capability.}
We adopt the compromised-prompt and compromised-observation capability class (the union of the SARC~\cite{sarc2026} and VeriGuard~\cite{veriguard2025} threat models). An attacker may (C1)~rewrite the system prompt, including operator-directive impersonation; (C2)~splice attacker-controlled text into every observation; and (C3)~persist both rewrites for the whole episode without reacting to LLM behavior (worst case for the defender). The attacker may \emph{not} modify the verifier, checkers, manager, simulator, or tool-layer validators, bypass \texttt{apply\_action}, or read the verifier's shadow state. Network, side-channel, supply-chain, and adaptive adversaries (tuning injections against observed verdicts) are out of scope (Sec.~\ref{sec:discussion}).

\paragraph{Failure semantics.}
Two primary outcomes and one derived. \emph{Material worsening}: the trajectory ends with strictly higher final penalty than the no-intervention default---worse than doing nothing. \emph{False recovery}: the trajectory reports terminal recovery while the live simulator state still contains violations. \emph{Attack success} (derived): the union of the two---any safety-relevant outcome the benign trajectory would have avoided. A recovery-rate drop without any of these is a liveness limitation, not a safety failure.

\paragraph{Instantiated attacks.}
We implement four attack families (detailed in Sec.~\ref{sec:eval:rq4}): \emph{prompt injection} (C1, a high-priority override pushing an aggressive fixed setpoint), \emph{observation poisoning} (C2, a fabricated ``all readings nominal'' stanza contradicting the violation list the verifier re-derives), \emph{stall} (C1, near-no-op setpoints that stay admissible while never recovering; a ``stall-RAG'' variant instructs the same), and \emph{magnitude redistribution} (C1/C2, driving a single retained branch toward failure while holding the overloaded-branch set and aggregate penalty within nominal bounds, masking per-branch overload behind nominal aggregate- and support-level readings).